\documentclass[12pt]{article}
\usepackage[utf8]{inputenc}
\usepackage[T1]{fontenc}
\usepackage[margin=1in]{geometry}
\usepackage{setspace}
\usepackage{xcolor}
\usepackage[colorlinks=true, urlcolor=blue!60!black, linkcolor=blue!60!black]{hyperref}
\setlength{\parindent}{0em}
\setlength{\parskip}{0.45em}
\renewcommand{\baselinestretch}{1.0}
\begin{document}
\begin{center}
{\Large \textit{Christian Legal Society v.\ Martinez}, 561 U.S.\ 661 (2010)\par}
\vspace{0.3em}
{\normalsize Case Brief --- One-Page Edition\par}
\end{center}

\textbf{Memory Jogger.} The all-comers case: student-group recognition as a limited public forum (\href{https://supreme.justia.com/cases/federal/us/561/661/}{561 U.S. 661}).

\textbf{Facts.} Hastings College of the Law gave recognized student groups funding, communication channels, and facilities, and the parties stipulated that it required every group to accept all comers as members and leaders. The Christian Legal Society chapter required officers and members to sign a Statement of Faith and was denied recognition.

\textbf{Procedural History.} Summary judgment for Hastings; the Ninth Circuit summarily affirmed.

\textbf{Votes.} Affirmed, 5--4. Ginsburg wrote for the Court; Stevens and Kennedy concurred; Alito dissented, joined by Roberts, Scalia, and Thomas.

\textbf{Holding.} Speech and association claims against a limited forum's membership rules are reviewed together under forum doctrine. A genuine all-comers policy is reasonable and viewpoint neutral, and \textit{Smith} bars a free-exercise exemption.

\textbf{Reasoning.} Hastings offered a subsidy rather than a prohibition, so \textit{Dale} and \textit{Jaycees} did not apply. The policy opens fee-funded leadership to all students, spares administrators from policing belief-versus-status lines, and encourages tolerance; it is ``textbook viewpoint neutral.'' Kennedy warned that infiltration or pretext would be a different case. Alito saw pretext and viewpoint discrimination under \textit{Healy v.\ James}.

\textbf{Examples.} \emph{Same-Side:} a uniform all-comers rule applied to a political club. \emph{Opposite-Side:} deregistering only a religious group for a faith requirement. \emph{Fence-Sitter:} opponents take over a small club.

\textbf{Critique.} Critics call the stipulated case artificial and the rule self-defeating for a diversity forum; progressives defend it as civic equality.

\textbf{Key Quotation.} ``The First Amendment shields CLS against state prohibition of the organization's expressive activity, however exclusionary that activity may be. But CLS enjoys no constitutional right to state subvention of its selectivity.'' --- Justice Ginsburg, Opinion of the Court, 561 U.S. at 669.

\end{document}
